\subsection{An extension homotopy for the angle-forgotten carriers}
\label{mapsIII:carrier-extension-section}

The literal flat quotient on the punctured cusp need not extend
continuously across the central fibre. Its homotopy class as a map to a
torus does extend. This distinction gives a second proof of the
angle-forgotten rank-two nullhomotopy and settles the one-copy circle
question left open in the working-edition homotopy-status remark.

\begin{lemma}[Suspending a torus-valued map that extends up to homotopy]
\label{mapsIII:carrier-extension-lemma}
Let $X$ be a compact Hausdorff space, and let $N\subset X$ be a compact
manifold with boundary which is a closed neighbourhood and satisfies
$\operatorname{int}_X(N)=N\setminus\partial N$.
Thus its frontier in $X$ is exactly its manifold boundary.
Let $r:N\to[0,\epsilon]$ be continuous, with
$W=r^{-1}(0)$, $\partial N=r^{-1}(\epsilon)$, and put
$U=N\setminus W$. Suppose $U$ is connected and locally path connected.
Let $\alpha:U\to T^d=\mathbb R^d/\mathbb Z^d$ be continuous.
Suppose there is a continuous $h:N\to T^d$ for which
$\alpha$ and $h|_U$ induce the same homomorphism to
$\pi_1(T^d)=\mathbb Z^d$. Fix $x_0\in W$, and choose $h(x_0)=0$.
For a based continuous map $c:(T^d,0)\to(Z,z_0)$ to a compact
Hausdorff based space, define
\[
 G(x)=
 \begin{cases}
 [r(x),c(\alpha(x))],&x\in U,\\
 *,&x\in W\cup(X\setminus\operatorname{int}N),
 \end{cases}
 \qquad
 \Sigma_r Z=
 \frac{[0,\epsilon]\times Z}
 {(\{0,\epsilon\}\times Z)\cup([0,\epsilon]\times\{z_0\})}.
\]
Then $G:X\to\Sigma_rZ$ is based nullhomotopic, by a homotopy
stationary on $(X\setminus\operatorname{int}N)\cup\{x_0\}$.
\end{lemma}

\begin{proof}
The difference $\alpha-h|_U:U\to T^d$ induces the zero homomorphism
on fundamental groups. The lifting criterion for the covering
$\mathbb R^d\to T^d$ therefore supplies a continuous function
$L:U\to\mathbb R^d$ with
\[
 \alpha(x)=h(x)+[L(x)].
\]
For $s\in[0,1]$, put
\[
 A_s(x)=h(x)+[(1-s)L(x)],\qquad
 G_s(x)=[r(x),c(A_s(x))]\quad(x\in U),
\]
and set $G_s=*$ on $W$ and outside $\operatorname{int}N$.
These formulas agree on the outer boundary. They are jointly continuous
at $W$ and at $\partial N$, even if $L$ is unbounded toward $W$:
every neighbourhood of the suspension basepoint contains the images of
$[0,\delta)\times Z$ and $(\epsilon-\delta,\epsilon]\times Z$
for some $\delta>0$, by compactness of $Z$. This estimate is uniform
in $s$ and in $A_s(x)$. Thus $G_s$ is a based homotopy from $G$
to the map $x\mapsto[r(x),c(h(x))]$ on $N$.

Now use that $h$ is defined on all of $N$. Set
\[
 H_s(x)=[(1-s)r(x)+s\epsilon,c(h(x))]
 \quad(x\in N),
 \qquad H_s(x)=*\quad(x\notin\operatorname{int}N).
\]
This is continuous by the closed pasting lemma. The height on
$\partial N$ is always $\epsilon$, so the formulas agree there.
At $x_0$ the second coordinate is $c(0)=z_0$, so $H_s(x_0)=*$
for every $s$. Finally $H_0=G_1$ and $H_1=*$.
Concatenating the two homotopies proves the assertion.
\end{proof}

\begin{proposition}[Both angle-forgotten one-copy carriers are nullhomotopic]
\label{mapsIII:angle-forgotten-one-copy-null}
For the actual cusp neighbourhood $N_0=f^{-1}(|t|\leq\epsilon)$ and
the original flat marking, the maps
\[
 \widehat q_3:X\longrightarrow S^3,
 \qquad
 \widehat q_\gamma:X\longrightarrow S^2
\]
defined by forgetting the base angle and taking the reduced suspension
are based nullhomotopic. The assertion concerns their original formulas
and basepoint $x_0\in W$. It also gives based nullhomotopies of every
positive smash power of either map.
\end{proposition}

\begin{proof}
Write
\[
 \Lambda=\mathbb Z\langle\widehat\gamma,\widehat u,
                    \widehat w,\widehat\delta\rangle,
 \quad
 K=\mathbb Z\langle\widehat w,\widehat\delta\rangle,
 \quad
 \bar\Lambda=\Lambda/K.
\]
The geometric fundamental-group calculation in
Proposition~\ref{mapsI:top:prop:geometric-pi1} proves the following
local information before taking any global quotient:
\[
 \pi_1(N_0)=\bar\Lambda,\qquad
 \pi_1(N_0\setminus W)=\Lambda\rtimes_{M_0}\mathbb Z,
 \qquad
 (\lambda,n)\longmapsto\bar\lambda.
\]
Take the radius $\epsilon$ strictly inside the punctured-overlap chart.
Because the fibration is a submersion over a small annulus about
$|t|=\epsilon$, a radial product collar identifies the fundamental
groups of this closed neighbourhood with those of the slightly larger
open neighbourhood used in that calculation, and identifies the
punctured inclusions as well.  This same collar makes $N_0$ a closed
codimension-zero neighbourhood in $X$, with
$\operatorname{int}_X(N_0)=N_0\setminus\partial N_0$ and frontier
$\partial N_0=f^{-1}(|t|=\epsilon)$, as required by the lemma.
Here the generator of the second factor is the zero-section meridian.
It bounds the actual holomorphic disc
$t\mapsto(t,1,1)$ in the toric chart, so there is no additional
translation in the displayed inclusion. Also
$(M_0-I)\Lambda=K$, as follows directly from
$(M_0-I)(a,b,c,d)=(0,0,b,-a)$.

The native map
$\alpha:N_0\setminus W\to\bar T^2=
(\bar\Lambda\otimes\mathbb R)/\bar\Lambda$
therefore induces precisely $(\lambda,n)\mapsto\bar\lambda$.
There is a map $h:N_0\to\bar T^2$ inducing the identity on
$\pi_1(N_0)=\bar\Lambda$. To see existence without any extra
classification assumption, take a CW model for the compact smooth
manifold $N_0$, map its one-cells with the specified integral winding
numbers, and extend over each two-cell because its attaching word
has trivial image in $\bar\Lambda$. All higher extensions exist
because the universal cover $\mathbb R^2$ of $\bar T^2$ is
contractible, so $\pi_k(\bar T^2)=0$ for $k\geq2$.
Compose with the CW comparison map to obtain $h$ on $N_0$.
Translation in the target arranges $h(x_0)=0$ and changes none of
the induced fundamental-group homomorphisms.

For $\widehat q_3$, apply Lemma~\ref{mapsIII:carrier-extension-lemma}
with $r=|t|$ and with
\[
 c:\bar T^2\longrightarrow
 \bar T^2/\bigl((S^1\times\{0\})\cup(\{0\}\times S^1)\bigr)
 \cong S^2.
\]
Its reduced suspension is exactly the original
$\widehat q_3(x)=[|t|,c(\alpha(x))]$ on the punctured cusp.
For $\widehat q_\gamma$, compose $\alpha$ and $h$ with the first
coordinate character $\bar T^2\to S^1$ and take $c$ to be the
identity of the based circle. The same lemma gives the nullhomotopy.
Both homotopies fix $x_0$, so their smash products with any based maps
are well defined and nullhomotopic.
\end{proof}

\paragraph{Scope and preservation of the other proof.}
This argument uses the proved local fundamental-group inclusion and does
not produce a pointwise extension of the original flat character across
$W$. In particular, it does not alter the nonextension result for that
literal character. The explicit three-framed torus and its embedded
nullbordism provide a distinct proof of $\widehat q_3\simeq*$ and
are retained. The retained-angle map $q_2$ and its nonzero
$S^4$-valued postcomposition have different targets; the present
homotopy does not apply to them because their base-angle coordinate has
nonzero winding on the meridian that is killed in $N_0$.
